% 09-disk-modes.tex -- a disk-mode surrogate, and a possible route to limb darkening.
% Input by mostlyharmless.tex.  No preamble here.
% This section is a proposal to be explored, not a settled part of the derivation.

\section{A disk-mode surrogate, and a possible route to limb darkening}
\label{sec:disk-modes}

This section is different from the ones before it: it derives nothing that the code
implements yet.  It records a proposal, the measurement that motivates it, and the reason
it may change how limb darkening is approached later.  Everything here is to be explored.

\subsection{The problem it addresses}
\label{sec:disk-modes-problem}

A parameter search does not need a faithful visibility everywhere.  It needs a surrogate
that can be evaluated at a \emph{sparse} set of baselines, many times, over many parameter
sets, and that is good enough to locate the interesting parts of parameter space --- after
which the numerical model is used to decide.  Percent-level fidelity is not the goal;
being wrong by a \emph{systematic} amount large enough to move the best-fit point is.

The affine model of the deformation family is the current surrogate.  Its structure is one
2$\times$2 stretch,
\begin{equation}
  V_\text{aff}(u,v) \;=\; V_\text{sph}\!\left(\Lambda^{T}(u,v)\right) ,
  \label{eq:affine-surrogate}
\end{equation}
where $V_\text{sph}$ is the closed-form visibility of the brightness pattern on the unit
sphere and $\Lambda$ is the projected stretch of the principal ellipsoid.  Its cost
structure is the attractive part: the projection onto the harmonics is paid once per
parameter set, and each frequency then costs a short Bessel sum.

\subsection{What the stretch keeps, and what it loses}
\label{sec:disk-modes-error}

Write the deformed surface as $r(n)\,n$ with $n$ a unit direction, and let $n_\perp$ be its
component in the sky plane.  The sky position of the emitting point is
\begin{equation}
  \mathbf{x}_\perp(n) \;=\; r(n)\, n_\perp ,
  \label{eq:sky-map}
\end{equation}
which is \emph{radial-nonlinear}: $r$ depends on the full direction, including its
component along the line of sight, and therefore on $|n_\perp|$ in a way no constant matrix
reproduces.  The affine model replaces \eqref{eq:sky-map} by its linear part
$\Lambda n_\perp$.

That replacement is not a second-order effect, and the difference is measurable.  Comparing
the surrogate against a direct sum over the same samples, with the projection, the
truncation and the normalisation all held fixed, gives a residual of $0.0016$ at the
working amplitudes and $0.0208$ with the amplitudes halved --- a ratio of $0.48$, i.e.\ the
error is \emph{first order} in the deformation $\epsilon$.  The same comparison with no
deformation at all returns $0.0000$ to the precision printed, which is what shows the
machinery itself is sound and the residual is the linearisation, not a defect in the
projection.

So the affine model is a first-order surrogate whose error is comparable to $\epsilon$
itself.  For a search that may be acceptable.  It is worth knowing, because it says the
surrogate cannot be improved by raising the harmonic truncation: the pattern's content above
$l = 2$ is already $O(\epsilon^2)$, far below the error of the map.

\subsection{Composing the stretch with a disk series}
\label{sec:disk-modes-composition}

Keep the stretch, and put a series \emph{inside} it.  Let $Z_M$ be a mode of the unit disk,
with $M$ collecting its two indices, and write
\begin{equation}
  V(u,v) \;\propto\; \sum_M c_M\,\widehat{Z}_M\!\left(\Lambda^{T}(u,v)\right) ,
  \label{eq:disk-mode-series}
\end{equation}
where $\widehat{Z}_M$ is the closed-form transform of the mode, derived below.

The reason the stretch belongs inside the series is geometrical.  The silhouette of the
deformed star is \emph{nearly} an ellipse, and $\Lambda$ is exactly the map that takes that
ellipse to a circle.  Composing the two means the basis is applied on the domain where its
transforms are elementary, and the residual that the stretch cannot absorb --- the departure
of \eqref{eq:sky-map} from linearity --- is what the higher modes carry.  The affine model
of \eqref{eq:affine-surrogate} is then the leading term of the series rather than a
different formalism, so nothing is thrown away by extending it.

The error stops being a fixed $O(\epsilon)$ and becomes a truncation in $M$, which is a
knob: more modes, smaller error, at a cost that is still a short sum per frequency.

\paragraph{How the coefficients are obtained.}  They are a projection, not a fit, and the
formula is the disk analogue of the spherical one.  Write $\tilde I(y) = I(\Lambda
y)\,|\det\Lambda|$ for the image pulled back to the unit disk, so that $\int\tilde
I\,d^2y$ is the total flux.  Since
\begin{equation}
  V(u) \;=\; \int I(x)\,e^{-2\pi i\,u\cdot x}\,d^2x
         \;=\; \int \tilde I(y)\,e^{-2\pi i\,u\cdot\Lambda y}\,d^2y ,
  \label{eq:disk-pullback}
\end{equation}
expanding $\tilde I$ on the disk with
$c_M = \frac{n+1}{\pi}\int_{|y|\le1}\tilde I(y)\overline{Z_M(y)}\,d^2y$ reproduces
\eqref{eq:disk-mode-series}, the weight being the Zernike normalisation
$\int Z_M\overline{Z_{M'}} = \frac{\pi}{n+1}\delta_{MM'}$.  Substituting $x = \Lambda y$
turns the integral into a quadrature over the samples that already exist,
\begin{equation}
  c_M \;=\; \frac{n+1}{\pi}\sum_i I_i\,w_i\,\overline{Z_M\!\left(\Lambda^{-1}x_i\right)} ,
  \label{eq:disk-coefficients}
\end{equation}
with $w_i$ the projected area the sample carries in
\eqref{eq:deformation-visibility}.  So the only new ingredient is the basis function
evaluated at the pulled-back sky position: no new sampler, no new geometry, and a cost of
$O(NM)$ per parameter set against $O(N(l+1)^2)$ for the spherical projection.

Two remarks on \eqref{eq:disk-coefficients}.  First, because the visibility is normalised at
the origin, only the \emph{ratios} of the coefficients matter: the factor $(n+1)/\pi$, the
$|\det\Lambda|$ and the absolute scaling of the basis all cancel, so the basis needs no
particular normalisation.  Second, $\Lambda$ enters twice --- in the frequency argument of
the series and in the pull-back --- and it is a \emph{preconditioner} rather than part of the
answer: the exact series does not depend on it, only the rate of convergence does.  A
$\Lambda$ that circularises the silhouette well is what makes a short series possible.

\subsection{The transform of a disk mode}
\label{sec:disk-modes-transform}

With polar coordinates $\rho,\phi$ on the disk and $\psi$ the baseline azimuth as fixed in
section~\ref{sec:conventions}, the plane-wave factor integrates to a Bessel function of the
azimuthal order:
\begin{equation}
  \widehat{Z}_n^m(q,\psi) \;=\; 2\pi\,(-i)^m e^{im\psi}
  \int_0^1 R_{n,m}(\rho)\, J_m(2\pi q\rho)\, \rho \, d\rho ,
  \label{eq:zernike-transform}
\end{equation}
the angular integral being the standard representation of $J_m$.  The Zernike radial
polynomials $R_{n,m}$ are those already introduced in section~\ref{sec:zernike} for the even
family, and they carry the classical identity
\begin{equation}
  \int_0^1 R_{n,m}(\rho)\, J_m(v\rho)\, \rho\, d\rho
  \;=\; (-1)^{(n-m)/2}\,\frac{J_{n+1}(v)}{v} ,
  \label{eq:zernike-identity}
\end{equation}
with $n-m$ even.  So \emph{one mode costs one ordinary Bessel function}, and a series of $M$
modes costs $M$ of them.  That is the same shape of result as the even family's finite
Bessel sum.

The identity is \textbf{certificate \vtag{22}}, and it is worth being clear about why it can
be certified so cleanly.  The integrand is a polynomial times a Bessel function, so
Gauss--Legendre quadrature evaluates the integral to machine precision: any residual is the
identity's own, not the quadrature's.  The check runs over every $(n,m)$ with $n \le 8$ and
$n-m$ even, against the library's own $R_{n,m}$, so it certifies that the convention of
section~\ref{sec:zernike} --- in particular the normalisation $R_{n,m}(1) = 1$ --- is the one
the identity holds for.  That is the point worth having: an identity of this kind is easy to
quote with the wrong normalisation, and the check fails immediately if the conventions
disagree.

\subsection{Why this may also be the limb-darkening route}
\label{sec:disk-modes-limb}

Limb darkening is a statement about the disk.  For a disk seen face-on, the cosine of the
angle between the surface normal and the line of sight is
\begin{equation}
  \mu = \sqrt{1-\rho^2} ,
  \label{eq:limb-mu}
\end{equation}
a function of the impact parameter \emph{alone}.  A limb-darkening law $I(\mu)$ is therefore
a radial profile on the disk, and an expansion of it in Zernike modes with $m = 0$ is the
natural representation --- indeed the classical visibility of a linearly limb-darkened disk
is a two-term combination of ordinary Bessel functions, one term being the uniform-disk
Airy pattern and the other the limb correction.

Contrast this with working on the sphere.  There $\mu = n\cdot\hat e_\text{obs}$ is a
degree-one factor, and re-expanding the product of that factor with the map's harmonics
requires Clebsch--Gordan, or Gaunt, coefficients.  That is why limb darkening has been
expected to bring Gaunt into the calculation.  On the disk the same physics needs no Gaunt
at all: \eqref{eq:limb-mu} is already in the variable the basis is built from.

If that holds up, the disk-mode surrogate and limb darkening are the same piece of
machinery, and building the surrogate now would be building the basis that Stage~7 needs,
with the affine model as its first term.

\subsection{Caveats that decide it}
\label{sec:disk-modes-caveats}

Three things could prevent this from being as clean as it looks.

\begin{enumerate}
  \item \textbf{The silhouette is not a circle.}  Its radius varies by $O(\epsilon)$ about
    its mean.  Either the domain is scaled to cover it, at the cost of giving up the exact
    orthogonality that makes \eqref{eq:zernike-identity} work, or the mismatch is accepted as
    a truncation error of the same order as the one being resolved.  The stretch removes the
    bulk of this, but not all of it.
  \item \textbf{$\mu$ is only radial to first order.}  On a deformed surface the normal is
    not along $r\,n$: it tilts by $O(\epsilon)$.  So \eqref{eq:limb-mu} is exact for a
    sphere and approximate for the bodies of interest, and the tilt is the term that will
    limit Stage~7 once the radial part is exact.
  \item \textbf{The frame is a choice.}  The series must be defined with respect to the same
    $\Lambda$ that the rest of the model uses, and $\Lambda$ is built from an ellipsoid that
    is itself an approximation.  A different but equally defensible $\Lambda$ changes the
    coefficients, though not the series' convergence.
\end{enumerate}

\subsection{What would settle it}
\label{sec:disk-modes-test}

The test is the one the affine model already faces: project the actual projected brightness
of a deformed surface onto the disk modes, form the series, and compare with a direct sum
over the same samples at a sparse set of frequencies.  The affine model sits at $0.0016$ at
the working amplitudes and scales like $\epsilon$.  A successful surrogate should (i) beat
that residual by a clear factor at a modest mode count $M$, (ii) have a residual that falls
as $M$ grows rather than plateauing --- plateauing would mean the residual is geometry the
basis cannot see, not truncation --- and (iii) keep the per-frequency cost small against the
per-parameter-set projection, which is the resource that is cheap.

If all three hold, the surrogate replaces the affine model, and Stage~7 becomes a question
of radial modes rather than of Gaunt coefficients.

\subsection{Register}
\label{sec:disk-modes-register}

\begin{center}
\begin{tabular}{l l l}
\hline
tag & what is checked & kind \\
\hline
\vtag{22} & the Zernike transform identity, against the library's $R_{n,m}$ & numerical \\
\hline
\end{tabular}
\end{center}

One certificate is all this section can claim, because nothing else here is implemented: the
identity is a fact about the basis, while the surrogate itself is a proposal.  The experiment
of section~\ref{sec:disk-modes-test} is what would produce the next one.
